
\documentclass{article}
\usepackage{colortbl}
\usepackage{makecell}
\usepackage{multirow}
\usepackage{supertabular}

\begin{document}

\newcounter{utterance}

\centering \large Interaction Transcript for game `cladder', experiment `full\_v1.5\_default', episode 6543 with llama{-}3.1{-}8b{-}instruct.
\vspace{24pt}

{ \footnotesize  \setcounter{utterance}{1}
\setlength{\tabcolsep}{0pt}
\begin{supertabular}{c@{$\;$}|p{.15\linewidth}@{}p{.15\linewidth}p{.15\linewidth}p{.15\linewidth}p{.15\linewidth}p{.15\linewidth}}
    \# & \multicolumn{2}{c}{Player} && \multicolumn{2}{c}{Game Master} \\
    \hline

    \theutterance \stepcounter{utterance}  
    & & & \multicolumn{4}{p{0.6\linewidth}}{
        \cellcolor[rgb]{0.9,0.9,0.9}{
            \makecell[{{p{\linewidth}}}]{
                \texttt{\tiny{[P1$\langle$GM]}}
                \texttt{You are an expert in causal inference. The following question is not a typical commonsense query, but rather a meticulously designed question created by a professor specializing in causal inference, intended to assess the students' mastery of the course content.} \\
\texttt{Imagine a self{-}contained, hypothetical world with only the following conditions, and without any unmentioned factors or causal relationships: Kwox has a direct effect on muvy and swoq. Swoq has a direct effect on kwoz. Muvy has a direct effect on kwoz. For those who are not kwox, the probability of kwoz is 34\%. For those who are kwox, the probability of kwoz is 55\%. Will kwox increase the chance of kwoz?} \\
\texttt{Start your answer with 'Yes' or 'No', followed by additional reasoning or evidence to support your explanation.} \\
            }
        }
    }
     \\ \\

    \theutterance \stepcounter{utterance}  
    & \multicolumn{4}{p{0.6\linewidth}}{
        \cellcolor[rgb]{0.561,0.737,0.561}{
            \makecell[{{p{\linewidth}}}]{
                \texttt{\tiny{[P1$\rangle$GM]}}
                \texttt{No, kwox will not increase the chance of kwoz.} \\
\\ 
\texttt{To determine the causal relationship between kwox and kwoz, we need to consider the potential confounding variables and the causal pathways between the variables. } \\
\\ 
\texttt{First, let's examine the direct effects: kwox has a direct effect on muvy and swoq, and muvy has a direct effect on kwoz. This suggests that kwox may influence kwoz indirectly through muvy. However, we also know that swoq has a direct effect on kwoz, which could be a confounding variable.} \\
\\ 
\texttt{To assess the causal relationship between kwox and kwoz, we can use the backdoor criterion, which states that a variable X causes Y if there is a path from X to Y that does not contain any arrows pointing back to X (i.e., no backdoor paths).} \\
\\ 
\texttt{In this case, we have a backdoor path from kwox to kwoz through swoq: kwox → muvy → kwoz is not a backdoor path because it contains an arrow pointing back to kwox (kwox → muvy). However, the path kwox → swoq → kwoz is a backdoor path because it does not contain any arrows pointing back to kwox.} \\
\\ 
\texttt{To determine the causal effect of kwox on kwoz, we need to compare the probability of kwoz among those who are kwox to the probability of kwoz among those who are not kwox, while controlling for the confounding variable swoq.} \\
\\ 
\texttt{Let P(kwoz|kwox) be the probability of kwoz among those who are kwox, and P(kwoz|not kwox) be the probability of kwoz among those who are not kwox. We are given that P(kwoz|kwox) = 55\% and P(kwoz|not kwox) = 34\%.} \\
\\ 
\texttt{However, we also know that the probability of kwoz among those who are not kwox but have a high value of swoq is likely to be higher than 34\%. Since kwox has a direct effect on swoq, those who are kwox are likely to have a higher value of swoq than those who are not kwox.} \\
\\ 
\texttt{Therefore, the observed difference in the probability of kwoz between those who are kwox and those who are not kwox may be due to the confounding variable swoq, rather than a causal effect of kwox on kwoz.} \\
\\ 
\texttt{To establish a causal relationship between kwox and kwoz, we would need to show that the difference in the probability of kwoz between those who are kwox and those who are not kwox remains after controlling for the confounding variable swoq. However, based on the information provided} \\
            }
        }
    }
    & & \\ \\

    \theutterance \stepcounter{utterance}  
    & & & \multicolumn{2}{p{0.3\linewidth}}{
        \cellcolor[rgb]{0.9,0.9,0.9}{
            \makecell[{{p{\linewidth}}}]{
                \texttt{\tiny{[GM$|$GM]}}
                \texttt{No} \\
            }
        }
    }
    & & \\ \\

    \theutterance \stepcounter{utterance}  
    & & & \multicolumn{2}{p{0.3\linewidth}}{
        \cellcolor[rgb]{0.9,0.9,0.9}{
            \makecell[{{p{\linewidth}}}]{
                \texttt{\tiny{[GM$|$GM]}}
                \texttt{yes} \\
            }
        }
    }
    & & \\ \\

    \theutterance \stepcounter{utterance}  
    & & & \multicolumn{2}{p{0.3\linewidth}}{
        \cellcolor[rgb]{0.9,0.9,0.9}{
            \makecell[{{p{\linewidth}}}]{
                \texttt{\tiny{[GM$|$GM]}}
                \texttt{game\_result = LOSE} \\
            }
        }
    }
    & & \\ \\

\end{supertabular}
}

\end{document}
